\chapter{Architettura generale e trust zones}

\section{Overview architetturale}

RAM-USB è un'architettura a microservizi n-tier in cui ogni componente ha una responsabilità specifica. 
Questa separazione riduce il rischio che la compromissione di un componente comprometta l'intero sistema.

RAM-USB ha dieci componenti interni, organizzati in tre gruppi funzionali: 
\begin{itemize}
    \item \textbf{Accesso e registrazione}: Entry-Hub, Security-Switch, Database-Vault, Network-Manager, Certificate-Authority, Headscale.
    \item \textbf{Backup e restore}: Storage-Service.
    \item \textbf{Osservabilità}: Mosquitto, Metrics-Collector, Grafana.
\end{itemize}

La \autoref{fig:architecture-container} mostra questi dieci componenti e il protocollo con cui interagiscono.

\begin{figure}[H]
    \centering
    \includegraphics[width=0.75\textwidth]{figures/01-architecture-container.pdf}
    \caption{Diagramma dei container di RAM-USB.}
    \label{fig:architecture-container}
\end{figure}

\subsection{Entry-Hub}

Entry-Hub è il gateway pubblico del sistema, ovvero l'unico esposto a chiunque su internet. Per questo 
il sistema non gli concede fiducia implicita. Il suo processo gira quindi in un container distroless 
\cite{distroless-github}, seguendo il principio di riduzione dell'attack surface raccomandato da 
NIST SP 800-190 \cite{nist-sp-800-190}. 
Senza shell, infatti, la maggior parte delle vulnerabilità critiche note nei container non è sfruttabile 
\cite{wager2024distroless}

Entry-Hub espone tre endpoint HTTPS: \texttt{POST /api/health} e \texttt{POST /api/register} sono
raggiungibili da chiunque su internet, con certificato firmato dalla CA pubblica Let's Encrypt
\cite{aas2019letsencrypt}. \texttt{POST /api/login} usa lo stesso certificato, ma ha un listener separato,
raggiungibile solo da utenti con accesso alla rete privata di cui parleremo in dettaglio nel
\autoref{chap:rete}.

Quando riceve una richiesta, Entry-Hub valida il corpo JSON in due fasi: in decodifica, impone un
limite di dimensione e rifiuta qualunque campo non atteso nel payload; poi valida i singoli
campi email, password e, solo in fase di registrazione, la chiave SSH \cite{ramusb-code-eh-f-04}. Se una di
queste verifiche fallisce, risponde HTTP 400 senza specificare quale sia il problema e senza inoltrare la
richiesta agli strati più interni del sistema. Se la validazione riesce, inoltra la richiesta a
Security-Switch via mTLS verificando che il certificato ricevuto appartenga davvero a un Security-Switch.
Attende la risposta del Security-Switch e la rilancia al client.

\subsection{Security-Switch}

Security-Switch è lo smistatore di richieste del sistema. Si occupa di ricevere le richieste da Entry-Hub e 
contattare i servizi interni. In questo modo la responsabilità di comunicare con i gestori del database e 
della rete non è in mano al punto di ingresso del sistema. 

Accetta connessioni in ingresso solo da Entry-Hub, autenticato via mTLS e, visto che il sistema tratta 
tutti i componenti come non fidati, ri-valida l'input in modo indipendente. 
Lo scopo di questa rivalidazione è l'applicazione del principio di defence in depth \cite{nist-sp-800-207}.

Se la validazione riesce, inoltra la richiesta a Database-Vault via mTLS. 

Nel caso di una richiesta di registrazione, chiede al Network-Manager di creare l'identità del nuovo utente 
sulla rete mesh e la credenziale (pre-auth key) con cui il suo dispositivo potrà unirvisi. 

Nel caso di una richiesta di login, chiede al Network-Manager di concedere un accesso temporaneo di 12 ore 
per quell'utente verso Storage-Service. 
Solo quando l'utente avrà ottenuto l'accesso allo Storage-Service, potrà contattarlo e gestire i propri backup.


\subsection{Database-Vault}


\subsection{Storage-Service}

\subsection{Network-Manager}

\subsection{Certificate-Authority}

\subsection{Headscale}

\subsection{Mosquitto}

\subsection{Metrics-Collector}

\subsection{Grafana}

\section{Modello di comunicazione inter-servizio}

\section{Trust zones e confini di sicurezza}
\label{sec:trust-zones}
